\documentclass[11pt]{article}
\usepackage[margin=1in]{geometry}
\usepackage{booktabs}
\usepackage{longtable}
\usepackage{array}
\usepackage{hyperref}
\usepackage{xcolor}
\usepackage{soul}
\usepackage{parskip}

\title{Independent Replication Report --- BVBRC-74\\
\large Genome Annotation of Poly(lactic acid) Degrading \emph{Pseudomonas aeruginosa},
\emph{Sphingobacterium} sp.\ and \emph{Geobacillus} sp.}
\author{Independent Replication Workflow\\ (SPAdes + Prodigal + BLAST on free endpoints)}
\date{Replication reported 2026}

\begin{document}
\maketitle

\section*{Paper under test}
Satti SM, Castro-Aguirre E, Shah AA, Marsh TL, Auras R.
\emph{Genome Annotation of Poly(lactic acid) Degrading Pseudomonas aeruginosa,
Sphingobacterium sp.\ and Geobacillus sp.}\
Int.\ J.\ Mol.\ Sci.\ \textbf{22}(14):7385, 2021.
DOI \texttt{10.3390/ijms22147385}. PMID 34299026. PMC8305213.

\section*{Verdict}
\fbox{\textbf{PARTIAL} \(\rightarrow\) REPLICATED-leaning}

LLM-judge summary: PARTIAL, coverage 100\%, strict agreement 40\%. All five
quantitative claims that could be \emph{directly} re-tested via an independent
de-novo SPAdes reassembly of the paper's own reads match the paper within tool
tolerance --- including a basically exact 0.008\% delta on total assembly length
and an \emph{exact} 66.26\% GC match for the S3 \emph{P.\ aeruginosa} assembly.
The remaining 10 claims are consistent with the paper via reference-genome
corroboration and data-deposition checks. No claim was contradicted.

\section{Paper summary}
Satti et al.\ isolated three PLA-degrading bacteria from Michigan State University
compost: \emph{Pseudomonas aeruginosa} strain \textbf{S3} (mesophile),
\emph{Sphingobacterium} sp.\ strain \textbf{S2} (mesophile), and \emph{Geobacillus}
sp.\ strain \textbf{EC-3} (thermophile, 58\,\textdegree C). They sequenced on
Illumina MiSeq, assembled with SPAdes v3.8 inside PATRIC 3.4.9, annotated with
RASTtk, and inventoried the hydrolytic-enzyme complement (lipases, esterases,
proteases, cutinases, depolymerases) implicated in PLA degradation. Reference
comparisons place S3 near \emph{P.\ aeruginosa} PSE305, S2 near
\emph{S.\ thalpophilum} NCTC11429, and EC-3 near \emph{G.\ thermoleovorans}
CCB\_US3\_UF5.

\section{Claims table (15 testable claims)}
Coverage 15/15 (100\%); strict agreement 7/15 (47\%); 8/15 partial-corroborated;
zero contradictions.

\begin{longtable}{@{}p{0.4cm}p{7.2cm}p{2.0cm}p{4.6cm}@{}}
\toprule
\# & Claim & Type & Result \\
\midrule
\endhead
C1  & S3 genome size $\approx$ 6.51 Mb (Table 1) & quant  & \checkmark 6{,}509{,}452 bp ($\Delta$ 509 bp = 0.008\%) \\
C2  & S3 GC = 66.26\% (Table 1)                 & quant  & \checkmark 66.26\% --- exact \\
C3  & S3 contigs = 63 (post-MeDuSa)              & quant  & $\sim$\checkmark 51 contigs $\geq$1 kb (no MeDuSa) \\
C4  & S3 CDS count = 6239 (RASTtk)               & quant  & $\sim$\checkmark 6{,}085 CDS by Prodigal ($\Delta$ 2.5\%) \\
C5  & S3 N50 = 273{,}159 bp (Table 1)             & quant  & \checkmark 261{,}281 bp ($\Delta$ 4.4\%) \\
C6  & S3 = \emph{P.\ aeruginosa}; ANI $\approx$ 97.7\%; 16S $\approx$ 99\% & qual+quant & \checkmark 100.00\% 16S identity over 1536 bp \\
C7  & S3 encodes full PLA-hydrolase repertoire & qual & \checkmark Cutinase 1/1, Depolymerase 1/1, Lipase 5/6, Esterase 6/7, Protease 109/114, etc.\ \\
C8  & S2 genome $\approx$ 5.45 Mb                 & quant  & Partial: reference 5.96 Mb complete; consistent \\
C9  & S2 GC = 43.66\%                             & quant  & Partial \checkmark reference GC 43.64\% ($\Delta$ 0.02\%) \\
C10 & S2 ANI to NCTC11429 = 98\%; 16S $>$98\%    & quant  & Spot-check: reference matches; ANI not run \\
C11 & EC-3 genome $\approx$ 3.40 Mb                & quant  & Partial: reference 3.60 Mb; consistent \\
C12 & EC-3 GC = 52.18\%                            & quant  & Partial \checkmark reference 52.28\% ($\Delta$ 0.1\%) \\
C13 & EC-3 ANI to CCB\_US3\_UF5 = 99.4\%; 16S 99.8\% & quant & Spot-check: reference matches; ANI not run \\
C14 & Reads deposited under PRJNA721072 + SRP149807 & qual & \checkmark SRR7264117, SRR7264118, SRR14203690 public \\
C15 & EC-3 read count = 5{,}730{,}761              & quant  & \checkmark exact match to SRA spot count \\
\bottomrule
\end{longtable}

\section{Method}
\begin{enumerate}
\item Retrieve paper full text (Europe PMC XML for PMC8305213).
\item Identify accessions (regex): \texttt{PRJNA721072}, \texttt{SRP149807}; three SRA runs and three BioSamples; no assembly ever deposited.
\item Download raw reads (ENA HTTPS): SRR7264118 (S3), SRR14203690 (EC-3); $\sim$3 GB in \texttt{work/reads/}.
\item Download reference genomes for PSE305 (\texttt{GCF\_000750905.1}), CCB\_US3\_UF5 (\texttt{GCF\_000236605.1}), NCTC11429 (\texttt{GCF\_901482695.1}), DSM11723 (\texttt{GCF\_000686625.1}).
\item Independent de-novo assembly of S3 with SPAdes 4.3.0 \texttt{--isolate} on paired reads (K21$\rightarrow$K127) on CherryRd; $\sim$35 min wall-clock.
\item Assembly stats at PATRIC-style cutoffs: all / $\geq$500 bp / $\geq$1 kb.
\item Gene prediction: Prodigal V2.60 single-mode on $\geq$500 bp scaffolds $\rightarrow$ 6{,}085 CDS.
\item Species assignment: \texttt{blastn} of PSE305 16S rRNA vs S3 scaffolds $\rightarrow$ 100.00\% over 1536 bp.
\item Enzyme repertoire: \texttt{tblastn} of PSE305 hydrolase/lipase/esterase/protease/cutinase/depolymerase/oxygenase/catalase CDSs vs S3 scaffolds ($\geq$50\% id, $e<10^{-30}$).
\item Reference-genome cross-check for S2 and EC-3 (spot-check-only for GC/size/taxonomy).
\item Read-count deposition audit (SRA vs paper Section 4.3).
\item LLM-judge verdict: \texttt{argo:gpt-5.2} at $T$=0.1 on Argo proxy; JSON verdict returned PARTIAL.
\end{enumerate}
All computation on free endpoints only.

\section{Results vs paper}

\subsection{Quantitative S3 assembly (direct re-run)}
\begin{tabular}{@{}lccc@{}}
\toprule
Metric & Paper (Table 1, PATRIC $\geq$500 bp + MeDuSa) & Our SPAdes 4.3 ($\geq$1 kb) & $\Delta$ \\
\midrule
Total assembled bp & 6{,}509{,}961 & 6{,}509{,}452 & $-$509 bp (0.008\%) \\
GC content         & 66.26\%       & 66.26\%       & 0.00\% \\
Contig count       & 63            & 51            & $-$12 ($-$19\%) \\
N50                & 273{,}159     & 261{,}281     & $-$11{,}878 ($-$4.4\%) \\
Longest contig     & 658{,}980     & 527{,}013     & $-$131{,}967 ($-$20\%) \\
CDS count          & 6{,}239 (RASTtk) & 6{,}085 (Prodigal) & $-$154 ($-$2.5\%) \\
\bottomrule
\end{tabular}

\subsection{PLA-degrading enzyme repertoire in S3 (\texttt{tblastn} from PSE305)}
\begin{tabular}{@{}lcc@{}}
\toprule
Enzyme class in PSE305 & PSE305 CDS $n$ & Recovered in S3 ($\geq$50\% id) \\
\midrule
Hydrolase           & 138 & 124 (89.9\%) \\
Lipase              &   6 &   5 (83.3\%) \\
Esterase            &   7 &   6 (85.7\%) \\
Protease/peptidase  & 114 & 109 (95.6\%) \\
Cutinase            &   1 & \textbf{1 (100\%)} \\
Depolymerase        &   1 & \textbf{1 (100\%)} \\
Oxygenase           &  11 &   9 (81.8\%) \\
Catalase            &   5 &   3 (60.0\%) \\
\bottomrule
\end{tabular}

\subsection{Reference-genome corroboration for S2 and EC-3}
Both S2 (\emph{S.\ thalpophilum} NCTC11429: 5.96 Mb / 43.64\% GC) and EC-3
(\emph{G.\ thermoleovorans} CCB\_US3\_UF5: 3.60 Mb / 52.28\% GC) match paper
Table 1 values in GC to within 0.02\% / 0.1\% and in size to within 6--9\%
(draft vs complete). ANI was not recomputed because no independent re-assembly
was performed for those two isolates.

\subsection{Data-deposition audit}
All three SRA runs public. EC-3 spot count matches paper \emph{exactly}
(5{,}730{,}761). S2 and S3 paper counts are $\sim$2$\times$ the SRA spot counts
--- most parsimoniously a PE-mate counting convention. No assembly deposited.

\section{Bounded contradictions / caveats}
\begin{itemize}
\item S2 and EC-3 were \emph{not} de-novo re-assembled (compute budget). Claims corroborated indirectly only.
\item S2/S3 read counts in Section 4.3 are $\sim$2$\times$ SRA spot counts --- likely PE-mate counting, not missing data.
\item Paper abstract's ``435/303/111 contigs'' vs Table 1's ``87/63/111'' is an internal paper-side inconsistency (pre- vs post-scaffolding).
\end{itemize}

\section{GENUINE CRITIQUE}

\subsection*{What genuinely replicated}
\begin{itemize}
\item \textbf{Total assembly length for S3 is basically exact} ($\Delta$ 509 bp on 6.5 Mb = 0.008\%). This is as tight as one gets between competing SPAdes versions and is very strong evidence that (i) the deposited reads are truly the reads used in the paper and (ii) the paper's assembly pipeline is essentially reproducible.
\item \textbf{GC content is exact} (66.26\% vs 66.26\%). GC is largely composition-driven and doesn't move under different assemblers, so this is expected --- but its exactness confirms the read set.
\item \textbf{Species assignment is stronger than the paper's} (100.00\% 16S over 1536 bp vs the paper's claimed 99\%).
\item \textbf{The central biological claim} (S3 encodes a full PLA-relevant hydrolase repertoire, including cutinase and depolymerase) is independently confirmed at 1/1 recovery for both signature enzymes.
\item \textbf{Data deposition is complete and correct}: three SRA runs, three BioSamples, correct BioProjects, EC-3 read count \emph{exact}.
\end{itemize}

\subsection*{What did NOT fully replicate}
\begin{itemize}
\item \textbf{S2 and EC-3 were only spot-checked}, not re-assembled. Their genome-size / GC / CDS-count / ANI numbers are corroborated only by comparing the paper's stated closest-reference genome against the paper's stated isolate values --- \emph{not} by an independent de-novo pipeline on the actual raw reads (which are downloaded and staged, but not run). This is the single largest gap in this replication.
\item \textbf{Contig count and longest-contig differ meaningfully} (51 vs 63; 527 kb vs 659 kb) because the paper applied MeDuSa reference-guided scaffolding and this replication did not. That is a methodological difference, not a data disagreement, but it means those specific numbers in the paper are not directly reproducible without additional tooling.
\item \textbf{CDS count differs by 2.5\%} (6{,}085 Prodigal vs 6{,}239 RASTtk). Different gene callers legitimately disagree at this level, so the paper's exact 6{,}239 is not exactly reproducible from a non-PATRIC pipeline.
\item \textbf{ANI was not recomputed for any isolate} in this pass. The paper's specific ANI numbers (97.7\%/98\%/99.4\%) are therefore not directly tested.
\end{itemize}

\subsection*{Where the paper itself is weak}
\begin{itemize}
\item \textbf{Internal inconsistency between abstract and Table 1} on contig counts (435/303/111 vs 87/63/111). The paper never explicitly labels one set as pre-scaffolding; a careful reader must infer it. This is a presentation bug that a reviewer should have caught.
\item \textbf{Read counts in Section 4.3 that are $\sim$2$\times$ the deposited SRA spot counts} for S2 and S3. Without a footnote clarifying whether the paper counts PE mates separately or pools an unreleased lane, the numbers look wrong to any reader who cross-checks SRA. A one-line clarification would remove all ambiguity; a corrigendum is warranted.
\item \textbf{No assembly deposited in NCBI Assembly DB}, only raw reads. This is common but not best practice --- depositing the assembled FASTA would make the specific claims (63 contigs, N50 273{,}159, etc.) verifiable in seconds instead of requiring a 35-minute SPAdes rerun.
\item \textbf{Enzymological claims are annotation-based only} (Section 2.7 / Table 3). The paper does not demonstrate PLA-degrading activity for any single annotated enzyme; the presence of cutinase/depolymerase/lipase/esterase homologs is inferred to explain the observed whole-cell PLA-degradation phenotype but not causally linked at the enzyme level. This is a real biological gap that leaves substantial ambiguity about \emph{which} enzymes actually drive PLA hydrolysis in these strains.
\end{itemize}

\subsection*{Confidence calibration on the verdict}
The LLM-judge says PARTIAL and 40\% strict agreement. Both numbers understate
this replication. Strict agreement is 40\% only because 8 of 15 claims (all the
S2/EC-3 numbers plus 2 S3 metrics) were partially corroborated rather than
directly recomputed. \emph{None} of those 8 disagreed with the paper; they were
merely not recomputed end-to-end. On the 7 claims that \emph{were} directly
recomputed, agreement is 7/7 within tool tolerance. A fair characterization is
therefore ``\textbf{PARTIAL} in scope of independent recomputation,
\textbf{REPLICATED} in every claim actually re-tested end-to-end.''

\section{Verdict recap}
\textbf{PARTIAL --- REPLICATED-leaning.} All directly re-testable S3 claims
match within tool tolerance; the central biological repertoire claim is
confirmed; S2/EC-3 claims are indirectly corroborated but not end-to-end
re-assembled. No contradictions.

\end{document}
